\begin{notechapter}
\chapter{Recurrent Networks and Backpropagation Through Time}\label{ch:rnn-bptt}
\chaptersubtitle{One set of weights applied at every time step, and the gradient that flows back through all of them}

\begin{chapterintro}
A sentence, a speech signal or a sensor recording is a sequence whose length changes from
one example to the next, and in which what happens at time $t$ depends on what happened
before. A recurrent neural network (RNN) processes such a sequence one element at a time,
with the same weights at every step, and carries a \emph{hidden state} from one step to
the next. Training it needs the gradient of a loss that depends on the weights through
every time step. Unrolling the recurrence turns the RNN into an ordinary computational
graph with shared weights, and backpropagation on that graph is called
\emph{backpropagation through time} (BPTT). You will learn
\begin{itemize}[nosep]
  \item the equations of the simple (Elman) RNN and what its hidden state represents;
  \item how unrolling turns the RNN into a computational graph, and the two different
        meanings of ``$\partial L/\partial h_t$'';
  \item the BPTT recursion, its proof, and the gradients of all parameters, including the
        passage from one weight to the whole matrix $W_{hh}$;
  \item a complete numerical example, checked with PyTorch;
  \item why the gradient contributions of distant time steps vanish or explode, with a
        precise statement in terms of singular values;
  \item gradient clipping and truncated BPTT.
\end{itemize}
\prerequisites{backpropagation and the three local rules (\cref{ch:backpropagation});
computational graphs and shared weights (\cref{ch:computational-graphs}); the softmax
output and the cross-entropy loss (\cref{ch:sigmoid-softmax}).}
\end{chapterintro}

\section{The recurrence}\label{rnn-bptt:sec:recurrence}

\paragraph{Why not an MLP.} An MLP takes an input of fixed size. One could pad every
sequence to a maximal length and feed the concatenation to an MLP, but this has two
defects: the number of weights grows with the maximal length, and a pattern learned at
position $3$ is not reused at position $30$, because the weights that see position $3$ are
different from those that see position $30$. The RNN solves both problems in the same way
as the convolutional layer of \cref{ch:convolutional-layers}: it shares its weights, here
across time instead of across space.

\begin{definition}[Simple recurrent network]\label{rnn-bptt:def:rnn}
Let $x_1,\dots,x_T\in\R^{n}$ be an input sequence and $h_0\in\R^m$ an initial state (usually
$h_0 = 0$). The \emph{simple} or \emph{Elman} RNN (Elman, \emph{Finding Structure in Time}, 1990) computes, for
$t = 1,\dots,T$,
\begin{align*}
  a_t &= W_{hh}h_{t-1} + W_{xh}x_t + b_h &&\text{(pre-activation, }\R^m\text{)}\\
  h_t &= \phi(a_t) &&\text{(hidden state, }\R^m\text{)}\\
  o_t &= W_{ho}h_t + b_o &&\text{(output scores, }\R^K\text{)}\\
  \hat y_t &= \softmax(o_t) &&\text{(predicted class probabilities)}
\end{align*}
with a coordinatewise activation $\phi$ (usually $\tanh$), $W_{hh}\in\R^{m\times m}$,
$W_{xh}\in\R^{m\times n}$, $b_h\in\R^m$, $W_{ho}\in\R^{K\times m}$ and $b_o\in\R^K$. For targets
$y_1,\dots,y_T\in\{1,\dots,K\}$, the loss is
\[
  L = \sum_{t=1}^T\ell_t,\qquad \ell_t = -\log\hat y_{t,y_t}.
\]
\end{definition}

The weights carry no index $t$: the same five arrays are used at every step, so the number
of parameters, $m^2 + mn + m + Km + K$, does not depend on $T$. Unrolling the recursion
shows what the hidden state is:
\[
  h_t = \phi\bigl(W_{hh}\,\phi(W_{hh}\,h_{t-2} + W_{xh}x_{t-1} + b_h) + W_{xh}x_t + b_h\bigr)
  = \cdots = F_t(x_1,\dots,x_t;\,h_0) .
\]
It is a function of the whole prefix $x_{1:t} = (x_1,\dots,x_t)$, computed incrementally: $h_t$ is a
\emph{summary of the past}, of fixed size $m$, however long the past is. Whether it is a
\emph{good} summary, one that keeps the information the task needs, is what training
decides; \cref{rnn-bptt:sec:vanishing} shows why distant information is hard to learn.

\paragraph{Task shapes.} The definition attaches a loss to every step
(\emph{many-to-many}, as in tagging every word of a sentence). If only the last output
matters, as in classifying a whole sequence, then $L = \ell_T$ (\emph{many-to-one}); if one
input is followed by a generated sequence, the outputs are fed back as the next inputs
(\emph{one-to-many}, \cref{ch:sequence-models}). The gradient computation below covers all
cases: a step without a loss simply has $\ell_t = 0$.

\begin{figure}[ht]
\centering
\begin{tikzpicture}[
  st/.style={draw, circle, minimum size=8mm, inner sep=0pt, fill=bookaccent!12},
  io/.style={draw, rectangle, rounded corners=2pt, minimum size=7mm, inner sep=1pt},
  ar/.style={-{Stealth[length=2mm]}, thick},
  lab/.style={font=\scriptsize, fill=white, inner sep=1pt}]
  % rolled
  \node[io] (x) at (0,0) {$x_t$};
  \node[st] (h) at (0,1.6) {$h_t$};
  \node[io] (o) at (0,3.1) {$o_t$};
  \node[io, fill=bookpitfall!10] (l) at (0,4.5) {$\ell_t$};
  \draw[ar] (x) -- node[lab, right]{$W_{xh}$} (h);
  \draw[ar] (h) -- node[lab, right]{$W_{ho}$} (o);
  \draw[ar] (o) -- (l);
  \draw[ar] (h.west) .. controls +(-1.3,0.7) and +(-1.3,-0.7) .. node[lab, left]{$W_{hh}$} (h.west);
  \node[font=\small] at (0,-0.9) {rolled};
  \node at (1.5,1.6) {\Large$=$};
  % unrolled
  \foreach \t/\xx in {1/4.4,2/6.4,3/8.4} {
    \node[io] (x\t) at (\xx,0) {$x_{\t}$};
    \node[st] (h\t) at (\xx,1.6) {$h_{\t}$};
    \node[io] (o\t) at (\xx,3.1) {$o_{\t}$};
    \node[io, fill=bookpitfall!10] (l\t) at (\xx,4.5) {$\ell_{\t}$};
    \draw[ar] (x\t) -- node[lab, right]{$W_{xh}$} (h\t);
    \draw[ar] (h\t) -- node[lab, right]{$W_{ho}$} (o\t);
    \draw[ar] (o\t) -- (l\t);
  }
  \node[st, fill=white] (h0) at (2.6,1.6) {$h_0$};
  \draw[ar] (h0) -- node[lab, above]{$W_{hh}$} (h1);
  \draw[ar] (h1) -- node[lab, above]{$W_{hh}$} (h2);
  \draw[ar] (h2) -- node[lab, above]{$W_{hh}$} (h3);
  \node (dots) at (9.9,1.6) {$\cdots$};
  \draw[ar] (h3) -- (dots);
  \node[io, fill=bookpitfall!25] (L) at (10.4,4.5) {$L$};
  \draw[ar, gray] (l3) -- (L);
  \draw[ar, gray] (l2.north) .. controls +(0.5,0.6) and +(-0.8,0.6) .. (L.north west);
  \draw[ar, gray] (l1.north) .. controls +(0.8,1.0) and +(-1.0,0.9) .. (L.north);
  \node[font=\small] at (6.4,-0.9) {unrolled for $T = 3$};
\end{tikzpicture}
\caption{The RNN of \cref{rnn-bptt:def:rnn}. Left: the loop on $h_t$ means that $h_{t-1}$ is an input of
the computation of $h_t$. Right: the same network unrolled in time. Every arrow labelled
$W_{hh}$ carries the \emph{same} matrix; the biases and the softmax are not drawn.}
\label{rnn-bptt:fig:unrolled}
\end{figure}

\section{The unrolled computational graph}\label{rnn-bptt:sec:graph}

For a given sequence, the unrolled network of \cref{rnn-bptt:fig:unrolled} is a
computational graph in the sense of \cref{computational-graphs:def:graph}: a directed
acyclic graph whose nodes are $h_0$, the $a_t$, $h_t$, $o_t$, $\ell_t$, the loss $L$ and the
parameters. Its only unusual feature is that each parameter feeds $T$ nodes. By
\cref{computational-graphs:sec:params}, the gradient with respect to a shared parameter is
the \emph{sum} of the contributions of the nodes that use it. Everything else is the
backpropagation of \cref{ch:backpropagation}, applied to a graph that is as deep as the
sequence is long.

We use the notation of \cref{backpropagation:sec:local}: for an intermediate quantity
$u$, $\bar u = \nabla_uL$ is the gradient of the total loss $L$ with respect to $u$, a vector of the
shape of $u$. We also write
\[
  \delta_t = \bar a_t = \nabla_{a_t}L\in\R^m
\]
for the \emph{error signal} at time $t$, as $\delta^{(l)}$ was the error signal of layer $l$ in
\cref{backpropagation:thm:backprop}.

\begin{pitfall}[Two meanings of $\partial L/\partial h_t$]
The symbol $\partial L/\partial h_t$ is used in the literature with two different meanings, and mixing them
is the most common source of wrong BPTT formulas.
\begin{itemize}[nosep]
  \item The \emph{partial} derivative, with every other variable of the formula held fixed.
        In $L = \sum_s\ell_s$, $h_t$ appears explicitly only in $\ell_t$ (through $o_t$), so this
        derivative is $W_{ho}^\top\bar o_t$. It ignores that $h_{t+1},h_{t+2},\dots$ are themselves functions of $h_t$.
  \item The \emph{total} derivative through the graph: the change of $L$ when $h_t$ changes and
        all later quantities are recomputed. It contains the effect on all future losses.
\end{itemize}
In this chapter $\bar h_t$ always denotes the total derivative, the one gradient descent needs.
The partial derivative is the ``direct'' or ``local'' part of it.
\end{pitfall}

\section{Backpropagation through time}\label{rnn-bptt:sec:bptt}

In the graph, $h_t$ has exactly two children: $o_t$ (through which it affects $\ell_t$) and
$a_{t+1}$ (through which it affects everything that comes later). For $t = T$ only the first
exists. This is the whole content of the following theorem.

\begin{theorem}[BPTT]\label{rnn-bptt:thm:bptt}
In the RNN of \cref{rnn-bptt:def:rnn}, let $\bar o_t = \hat y_t - e_{y_t}$, where $e_k$ is the $k$-th standard basis
vector of $\R^K$. The gradients of $L$ satisfy the backward recursion
\begin{align*}
  \bar h_T &= W_{ho}^\top\bar o_T,\\
  \bar h_t &= W_{ho}^\top\bar o_t + W_{hh}^\top\delta_{t+1}\qquad(t = T-1,\dots,1),\\
  \delta_t &= \bar h_t\odot\phi'(a_t)\qquad(t = T,\dots,1),
\end{align*}
and the parameter gradients are sums over time:
\begin{gather*}
  \nabla_{W_{hh}}L = \sum_{t=1}^T\delta_t\,h_{t-1}^\top,\qquad
  \nabla_{W_{xh}}L = \sum_{t=1}^T\delta_t\,x_t^\top,\qquad
  \nabla_{b_h}L = \sum_{t=1}^T\delta_t,\\
  \nabla_{W_{ho}}L = \sum_{t=1}^T\bar o_t\,h_t^\top,\qquad
  \nabla_{b_o}L = \sum_{t=1}^T\bar o_t .
\end{gather*}
\end{theorem}

\begin{proof}
\emph{Outputs.} $o_t$ influences $L$ only through $\ell_t$, and the softmax with the cross-entropy
loss gives $\nabla_{o_t}\ell_t = \hat y_t - e_{y_t}$ (\cref{backpropagation:sec:output}).

\emph{Hidden states.} By the path-sum recursion of \cref{computational-graphs:lem:paths}
(the chain rule summed over the children of a node), the gradient at $h_t$ is the sum, over
its children, of the gradient at the child pulled back by the transposed local Jacobian.
The child $o_t = W_{ho}h_t + b_o$ contributes $W_{ho}^\top\bar o_t$ and the child $a_{t+1} = W_{hh}h_t + \cdots$ contributes
$W_{hh}^\top\bar a_{t+1}$ (rule through an affine map, \cref{backpropagation:sec:affine}). For $t = T$ there is no
child $a_{T+1}$.

\emph{Pre-activations.} $a_t$ has the single child $h_t = \phi(a_t)$, and the rule through a
coordinatewise activation (\cref{backpropagation:sec:activation}) gives
$\bar a_t = \bar h_t\odot\phi'(a_t)$.

\emph{Parameters.} $W_{hh}$ is used by the $T$ nodes $a_1,\dots,a_T$. For the node $a_t$ alone, with $h_{t-1}$
held fixed, the affine rule gives the contribution $\bar a_t\,h_{t-1}^\top$; summing over the nodes that
share the parameter (\cref{computational-graphs:sec:params}) gives $\nabla_{W_{hh}}L$. The other four
formulas are obtained in the same way: $W_{xh}$ and $b_h$ are used by the $a_t$, $W_{ho}$ and $b_o$ by
the $o_t$.
\end{proof}

The recursion runs backward in time, hence the name. The forward pass must store every
$h_t$ (and $a_t$, or $\phi'(a_t)$), since the backward pass needs them all.

\begin{figure}[ht]
\centering
\begin{tikzpicture}[
  nd/.style={draw, rounded corners=2pt, minimum width=10mm, minimum height=7mm, inner sep=1pt},
  ar/.style={-{Stealth[length=2mm]}, thick},
  bk/.style={-{Stealth[length=2mm]}, thick, bookpitfall, dashed},
  lab/.style={font=\scriptsize, inner sep=1pt}]
  \node[nd] (hm) at (0,0) {$h_{t-1}$};
  \node[nd] (at) at (2.2,0) {$a_t$};
  \node[nd, fill=bookaccent!15] (ht) at (4.4,0) {$h_t$};
  \node[nd] (at1) at (7.4,0) {$a_{t+1}$};
  \node[nd] (ht1) at (9.6,0) {$h_{t+1}$};
  \node[nd] (ot) at (4.4,2.0) {$o_t$};
  \node[nd] (lt) at (4.4,3.6) {$\ell_t$};
  \node (fut) at (11.4,0) {$\cdots$};
  \draw[ar] (hm) -- (at); \draw[ar] (at) -- node[lab, below]{$\phi$} (ht);
  \draw[ar] (ht) -- node[lab, below]{$W_{hh}$} (at1); \draw[ar] (at1) -- node[lab, below]{$\phi$} (ht1);
  \draw[ar] (ht1) -- (fut);
  \draw[ar] (ht) -- node[lab, left]{$W_{ho}$} (ot); \draw[ar] (ot) -- (lt);
  \draw[bk] (ot.east) to[bend left=40] node[lab, right, text=bookpitfall]{$W_{ho}^\top\bar o_t$} (ht.north east);
  \draw[bk] (at1.south west) to[bend left=35] node[lab, below, text=bookpitfall]{$W_{hh}^\top\delta_{t+1}$} (ht.south east);
  \draw[bk] (ht.north west) to[bend right=35] node[lab, above, text=bookpitfall]{$\delta_t = \bar h_t\odot\phi'(a_t)$} (at.north east);
\end{tikzpicture}
\caption{The two children of $h_t$ in the unrolled graph. Forward computations in black; the
backward pass (dashed) adds the gradients coming back from $o_t$ (the current loss) and from
$a_{t+1}$ (all future losses), then multiplies coordinatewise by $\phi'(a_t)$.}
\label{rnn-bptt:fig:children}
\end{figure}

\paragraph{The recursion unrolled.} Iterating the recursion shows where each part of the
gradient comes from. Write $J_r = \partial h_r/\partial h_{r-1} = \diag(\phi'(a_r))\,W_{hh}$ for the Jacobian of one step,
and, for $s\geq t$,
\[
  \frac{\partial h_s}{\partial h_t} = J_sJ_{s-1}\cdots J_{t+1}\qquad(\text{the identity for }s = t).
\]

\begin{proposition}[Contributions of the future losses]\label{rnn-bptt:prop:unrolled}
\[
  \bar h_t = \sum_{s=t}^T\Bigl(\frac{\partial h_s}{\partial h_t}\Bigr)^{\!\top}W_{ho}^\top\,\bar o_s .
\]
\end{proposition}

\begin{proof}
By backward induction on $t$. For $t = T$ the sum has the single term $W_{ho}^\top\bar o_T$. If the formula holds at $t+1$,
then $W_{hh}^\top\delta_{t+1} = W_{hh}^\top\diag(\phi'(a_{t+1}))\bar h_{t+1} = J_{t+1}^\top\bar h_{t+1}$, and by the induction hypothesis
\[
  \bar h_t = W_{ho}^\top\bar o_t + J_{t+1}^\top\sum_{s=t+1}^T\Bigl(\frac{\partial h_s}{\partial h_{t+1}}\Bigr)^{\!\top}W_{ho}^\top\bar o_s
  = W_{ho}^\top\bar o_t + \sum_{s=t+1}^T\Bigl(\frac{\partial h_s}{\partial h_{t+1}}J_{t+1}\Bigr)^{\!\top}W_{ho}^\top\bar o_s ,
\]
and $\frac{\partial h_s}{\partial h_{t+1}}J_{t+1} = \frac{\partial h_s}{\partial h_t}$.
\end{proof}

The loss at time $s$ reaches the state at time $t$ through a product of $s - t$ Jacobians. This
product is the object of \cref{rnn-bptt:sec:vanishing}.

\section{From one weight to the matrix \texorpdfstring{$W_{hh}$}{W\_hh}}\label{rnn-bptt:sec:weight}

The formula $\nabla_{W_{hh}}L = \sum_t\delta_th_{t-1}^\top$ hides two steps that are worth doing once by hand: the
derivative with respect to a matrix, and the sum over time. We follow one weight, then
assemble the matrix.

\paragraph{One weight.} The weight $(W_{hh})_{kl}$ connects unit $l$ of $h_{t-1}$ to unit $k$ of $a_t$, at every $t$. It
influences $L$ through the $mT$ numbers $a_{t,i}$, so by the chain rule
\begin{equation}\label{rnn-bptt:eq:one-weight}
  \frac{\partial L}{\partial(W_{hh})_{kl}} = \sum_{t=1}^T\sum_{i=1}^m\frac{\partial L}{\partial a_{t,i}}\,
  \frac{\partial a_{t,i}}{\partial(W_{hh})_{kl}} .
\end{equation}
Here $\partial L/\partial a_{t,i} = \delta_{t,i}$ is a total derivative (the gradient at node $a_t$), while
$\partial a_{t,i}/\partial(W_{hh})_{kl}$ is the local derivative of the formula
$a_{t,i} = \sum_j(W_{hh})_{ij}(h_{t-1})_j + \cdots$, with $h_{t-1}$ held fixed:
\[
  \frac{\partial a_{t,i}}{\partial(W_{hh})_{kl}} = \begin{cases}(h_{t-1})_l & \text{if } i = k,\\ 0 & \text{if } i\neq k,\end{cases}
  \qquad\text{that is}\qquad \frac{\partial a_{t,i}}{\partial(W_{hh})_{kl}} = [i = k]\,(h_{t-1})_l .
\]
The indicator collapses the sum over $i$ in \cref{rnn-bptt:eq:one-weight}:
\begin{equation}\label{rnn-bptt:eq:entry}
  \frac{\partial L}{\partial(W_{hh})_{kl}} = \sum_{t=1}^T\delta_{t,k}\,(h_{t-1})_l .
\end{equation}
For each $t$: the error signal at the head of the connection, times the activity at its tail.

\paragraph{The three-index array.} For a fixed $t$, the $m^3$ local derivatives
$\partial a_{t,i}/\partial(W_{hh})_{kl}$ form a three-index array $\partial a_t/\partial W_{hh}$. Fixing the output index $i$ gives an
$m\times m$ slice, the derivative of the single number $a_{t,i}$ with respect to the whole matrix: it
is zero except in row $k = i$, which contains $h_{t-1}^\top$ (\cref{rnn-bptt:fig:slices}). The array is very sparse,
with $m^2$ non-zero entries out of $m^3$, because $a_{t,i}$ depends only on row $i$ of $W_{hh}$.

\begin{figure}[ht]
\centering
\begin{tikzpicture}[scale=1.1]
  \node[inner sep=0pt, opacity=0.55] (back) at (6.6,5.0) {%
    $\begin{bNiceMatrix}
      0 & 0 & \Cdots & 0 \\
      \Vdots & \Vdots & & \Vdots \\
      0 & 0 & \Cdots & 0 \\
      (h_{t-1})_1 & (h_{t-1})_2 & \Cdots & (h_{t-1})_m
    \end{bNiceMatrix}$};
  \node[font=\scriptsize, text=bookaccent] at ($(back.east)+(0.45,0.9)$) {$i = m$};
  \node[inner sep=0pt, opacity=0.8] (middle) at (3.3,2.5) {%
    $\begin{bNiceMatrix}
      0 & 0 & \Cdots & 0 \\
      \Vdots & \Vdots & & \Vdots \\
      (h_{t-1})_1 & (h_{t-1})_2 & \Cdots & (h_{t-1})_m \\
      \Vdots & \Vdots & & \Vdots \\
      0 & 0 & \Cdots & 0
    \end{bNiceMatrix}$};
  \node[font=\scriptsize, text=bookaccent] at ($(middle.east)+(0.45,1.1)$) {$i = k$};
  \node[inner sep=0pt, fill=white] (front) at (0,0) {%
    $\begin{bNiceMatrix}
      (h_{t-1})_1 & (h_{t-1})_2 & \Cdots & (h_{t-1})_m \\
      0 & 0 & \Cdots & 0 \\
      \Vdots & \Vdots & & \Vdots \\
      0 & 0 & \Cdots & 0
    \end{bNiceMatrix}$};
  \node[font=\scriptsize, text=bookaccent] at ($(front.east)+(0.4,0.9)$) {$i = 1$};
  \draw[dotted, gray] (back.north west) -- (front.north west);
  \draw[dotted, gray] (back.north east) -- (front.north east);
  \draw[dotted, gray] (back.south east) -- (front.south east);
  \draw[->, thick, bookpitfall!80] ($(front.south west)+(0.1,-0.25)$) --
    node[midway, below, font=\scriptsize]{$l$ (column of $W_{hh}$)} ($(front.south east)+(-0.1,-0.25)$);
  \draw[->, thick, bookexample!80!black] ($(front.north west)+(-0.25,-0.1)$) --
    node[midway, left, font=\scriptsize, align=right]{$k$ (row\\ of $W_{hh}$)} ($(front.south west)+(-0.25,0.1)$);
  \draw[->, thick, bookaccent] ($(front.east)+(1.0,-0.4)$) --
    node[midway, below right, font=\scriptsize]{$i$ (slices)} ++(6.6,5.0);
\end{tikzpicture}
\caption{The array $\partial a_t/\partial W_{hh}$, drawn as $m$ slices of size $m\times m$. Slice $i$ is
$\partial a_{t,i}/\partial W_{hh} = e_i\,h_{t-1}^\top$: only its row $i$ is non-zero, and it contains $h_{t-1}^\top$.}
\label{rnn-bptt:fig:slices}
\end{figure}

\paragraph{Contraction.} The chain rule \cref{rnn-bptt:eq:one-weight} for all $(k,l)$ at once multiplies
the $i$-th slice by the number $\delta_{t,i}$ and adds the slices:
\[
  \sum_{i=1}^m\delta_{t,i}\,\frac{\partial a_{t,i}}{\partial W_{hh}} = \sum_{i=1}^m\delta_{t,i}\,e_i\,h_{t-1}^\top
  = \Bigl(\sum_{i=1}^m\delta_{t,i}e_i\Bigr)h_{t-1}^\top = \delta_t\,h_{t-1}^\top .
\]
The contraction of the sparse array with the error signal is an outer product, which is
the matrix of the entries \cref{rnn-bptt:eq:entry} at time $t$. This is the computation of
\cref{backpropagation:sec:affine}, where $\partial y/\partial W$ was the same array for a single affine map.

\paragraph{Sum over time.} The only new ingredient of the recurrent case is the sum over
$t$: $W_{hh}$ is used at every step, so
\[
  \nabla_{W_{hh}}L = \sum_{t=1}^T\delta_t\,h_{t-1}^\top
  = \begin{bmatrix}\delta_1 & \delta_2 & \cdots & \delta_T\end{bmatrix}
    \begin{bmatrix}h_0^\top\\ h_1^\top\\ \vdots\\ h_{T-1}^\top\end{bmatrix} = \Delta H_{-}^\top ,
\]
where $\Delta\in\R^{m\times T}$ stores the error signals as columns and $H_-\in\R^{m\times T}$ the states $h_0,\dots,h_{T-1}$. The
second form is how the sum is computed in practice: one matrix product after the backward
pass, instead of $T$ outer products.

\begin{keyidea}[The gradient of a shared weight]
Each use of $W_{hh}$ contributes \emph{(error signal at the output of the use)} $\times$ \emph{(input of the use)}$^\top$,
computed as if that use were the only one, and the contributions are added. The coupling
between time steps is entirely contained in the error signals $\delta_t$, which carry the effect of
all the future losses.
\end{keyidea}

\section{The algorithm}\label{rnn-bptt:sec:algorithm}

\begin{algorithm}[ht]
\caption{Backpropagation through time for one sequence}\label{rnn-bptt:alg:bptt}
\begin{algorithmic}[1]
\For{$t = 1,\dots,T$} \Comment{forward pass: store $h_t$ and $\phi'(a_t)$}
  \State $a_t\gets W_{hh}h_{t-1} + W_{xh}x_t + b_h$;\quad $h_t\gets\phi(a_t)$;\quad $\hat y_t\gets\softmax(W_{ho}h_t + b_o)$
\EndFor
\State $g\gets 0\in\R^m$ \Comment{$g = W_{hh}^\top\delta_{t+1}$, the gradient arriving from the future}
\For{$t = T,\dots,1$} \Comment{backward pass}
  \State $\bar o_t\gets\hat y_t - e_{y_t}$;\quad $\bar h_t\gets W_{ho}^\top\bar o_t + g$;\quad $\delta_t\gets\bar h_t\odot\phi'(a_t)$
  \State accumulate $\nabla_{W_{ho}}\mathrel{+}=\bar o_th_t^\top$, $\nabla_{b_o}\mathrel{+}=\bar o_t$, $\nabla_{W_{hh}}\mathrel{+}=\delta_th_{t-1}^\top$, $\nabla_{W_{xh}}\mathrel{+}=\delta_tx_t^\top$, $\nabla_{b_h}\mathrel{+}=\delta_t$
  \State $g\gets W_{hh}^\top\delta_t$
\EndFor
\end{algorithmic}
\end{algorithm}

\paragraph{Cost.} Both passes do $O(T(m^2 + mn + Km))$ operations, as for an MLP of depth $T$. The memory is
$O(Tm)$ for the stored states, which is what limits the length of the sequences in practice and
motivates truncation (\cref{rnn-bptt:sec:truncated}). Unlike the layers of an MLP, the time steps
cannot be computed in parallel: $h_t$ needs $h_{t-1}$.

\section{A worked example}\label{rnn-bptt:sec:example}

\begin{example}[BPTT with $m = 2$, $T = 3$]\label{rnn-bptt:ex:numbers}
Take $\phi = \tanh$, a scalar input ($n = 1$), $K = 2$ classes, $h_0 = 0$, zero biases and
\[
  W_{hh} = \begin{bmatrix}0.5 & -0.3\\ 0.2 & 0.4\end{bmatrix},\qquad
  W_{xh} = \begin{bmatrix}1\\ -0.5\end{bmatrix},\qquad
  W_{ho} = \begin{bmatrix}1 & -1\\ 0.5 & 0.5\end{bmatrix},
\]
with inputs $x = (1,\ 0.5,\ -1)$ and targets $y = (1,\ 2,\ 1)$. All numbers are rounded to four digits.

\emph{Forward pass.}
\begin{center}
\begin{tabular}{c|cccc}
$t$ & $a_t$ & $h_t$ & $\hat y_t$ & $\ell_t$\\\hline
$1$ & $(1,\ -0.5)$ & $(0.7616,\ -0.4621)$ & $(0.7454,\ 0.2546)$ & $0.2939$\\
$2$ & $(1.0194,\ -0.2825)$ & $(0.7696,\ -0.2752)$ & $(0.6895,\ 0.3105)$ & $1.1696$\\
$3$ & $(-0.5326,\ 0.5438)$ & $(-0.4874,\ 0.4959)$ & $(0.2714,\ 0.7286)$ & $1.3041$
\end{tabular}
\end{center}
For instance $a_2 = W_{hh}h_1 + W_{xh}x_2 = (0.5\cdot0.7616 + 0.3\cdot0.4621 + 0.5,\ 0.2\cdot0.7616 - 0.4\cdot0.4621 - 0.25) = (1.0194,\ -0.2825)$. The
loss is $L = 2.7675$.

\emph{Backward pass}, from $t = 3$, with $\phi'(a_t) = 1 - h_t^2$:
\begin{center}
\begin{tabular}{c|ccc}
$t$ & $\bar o_t = \hat y_t - e_{y_t}$ & $\bar h_t = W_{ho}^\top\bar o_t + W_{hh}^\top\delta_{t+1}$ & $\delta_t = \bar h_t\odot(1 - h_t^2)$\\\hline
$3$ & $(-0.7286,\ 0.7286)$ & $(-0.3643,\ 1.0929)$ & $(-0.2778,\ 0.8242)$\\
$2$ & $(0.6895,\ -0.6895)$ & $(0.3707,\ -0.6212)$ & $(0.1511,\ -0.5742)$\\
$1$ & $(-0.2546,\ 0.2546)$ & $(-0.1666,\ 0.1070)$ & $(-0.0700,\ 0.0841)$
\end{tabular}
\end{center}
At $t = 2$, the two children contribute $W_{ho}^\top\bar o_2 = (0.3448,\ -1.0343)$ (current loss) and
$W_{hh}^\top\delta_3 = (0.0259,\ 0.4130)$ (future loss). The second is not small: the second
coordinate of $h_2$ matters more through $\ell_3$ than through $\ell_2$.

\emph{Gradient of $W_{hh}$.} The contributions $\delta_th_{t-1}^\top$ are
\[
  \delta_1h_0^\top = 0,\qquad
  \delta_2h_1^\top = \begin{bmatrix}0.1151 & -0.0698\\ -0.4373 & 0.2653\end{bmatrix},\qquad
  \delta_3h_2^\top = \begin{bmatrix}-0.2138 & 0.0765\\ 0.6343 & -0.2268\end{bmatrix},
\]
and their sum is
\[
  \nabla_{W_{hh}}L = \begin{bmatrix}-0.0987 & 0.0066\\ 0.1970 & 0.0385\end{bmatrix}.
\]
The first contribution vanishes because $h_0 = 0$: at the first step, $W_{hh}$ multiplies the zero vector.
The other gradients are $\nabla_{W_{xh}}L = \sum_t\delta_tx_t = (0.2834,\ -1.0271)$, $\nabla_{b_h}L = (-0.1966,\ 0.3341)$,
$\nabla_{b_o}L = (-0.2938,\ 0.2938)$ and
\[
  \nabla_{W_{ho}}L = \begin{bmatrix}0.6918 & -0.4334\\ -0.6918 & 0.4334\end{bmatrix}.
\]
The two rows of $\nabla_{W_{ho}}L$ and the two entries of $\nabla_{b_o}L$ are opposite because $\hat y_{t,1} + \hat y_{t,2} = 1$, so the two
coordinates of every $\bar o_t$ are opposite.
\end{example}

Automatic differentiation computes the same numbers from the forward pass alone, which is
how every gradient of this chapter was checked:
\begin{lstlisting}[language=Python]
import torch
torch.set_default_dtype(torch.float64)
Whh = torch.tensor([[0.5, -0.3], [0.2, 0.4]], requires_grad=True)
Wxh = torch.tensor([[1.0], [-0.5]], requires_grad=True)
Who = torch.tensor([[1.0, -1.0], [0.5, 0.5]], requires_grad=True)
h, L = torch.zeros(2), 0.0
for x, y in zip([1.0, 0.5, -1.0], [0, 1, 0]):          # classes numbered from 0
    h = torch.tanh(Whh @ h + Wxh @ torch.tensor([x]))
    L = L - torch.log_softmax(Who @ h, dim=0)[y]
L.backward()
print(L.item(), Whh.grad)   # 2.7675  [[-0.0987, 0.0066], [0.1970, 0.0385]]
\end{lstlisting}

\section{Vanishing and exploding gradients}\label{rnn-bptt:sec:vanishing}

By \cref{rnn-bptt:prop:unrolled}, the loss at time $s$ contributes to $\bar h_t$ through the product
\[
  \frac{\partial h_s}{\partial h_t} = \prod_{r=s}^{t+1}\diag\bigl(\phi'(a_r)\bigr)\,W_{hh}\qquad(\text{factors ordered from }r = s\text{ on the left to }r = t+1)
\]
of $k = s - t$ Jacobians. Products of many matrices tend to shrink or grow exponentially in $k$,
which was identified as the central difficulty of training RNNs by Hochreiter (diploma
thesis, 1991) and by Bengio, Simard and Frasconi (\emph{Learning Long-Term Dependencies
with Gradient Descent is Difficult}, 1994).

\begin{proposition}[A bound on the Jacobian product]\label{rnn-bptt:prop:bound}
Let $\gamma = \sup_a\abs{\phi'(a)}$ ($\gamma = 1$ for $\tanh$, $\gamma = \tfrac14$ for the logistic sigmoid) and let $s_1 = \norm{W_{hh}}_2$
be the largest singular value of $W_{hh}$. Then, for $s\geq t$,
\[
  \Bigl\|\frac{\partial h_s}{\partial h_t}\Bigr\|_2\leq(\gamma s_1)^{s-t} .
\]
\end{proposition}

\begin{proof}
The spectral norm is submultiplicative, and the norm of a diagonal matrix is the largest
absolute value of its entries, so $\norm{\diag(\phi'(a_r))W_{hh}}_2\leq\gamma s_1$ for every factor.
\end{proof}

Two consequences follow, which are the precise form of the statements of
Pascanu, Mikolov and Bengio (\emph{On the Difficulty of Training Recurrent Neural Networks}, 2013).
\begin{itemize}
  \item \textbf{Vanishing.} If $\gamma s_1 < 1$, the contribution of the loss at time $s$ to $\bar h_t$ decreases at least
        geometrically with the distance $s - t$, whatever the inputs. The gradient then
        contains almost no information about long-range dependencies, and the network
        cannot learn them by gradient descent.
  \item \textbf{Exploding.} The contributions can grow exponentially only if $\gamma s_1 > 1$: this
        condition is \emph{necessary} for explosion, but not sufficient, because the factors
        $\phi'(a_r)$ can be much smaller than $\gamma$ when the units saturate.
\end{itemize}
Pascanu et al.\ state these conditions with the largest \emph{eigenvalue modulus} of $W_{hh}$ (its
spectral radius $\rho$). For a symmetric matrix, or more generally a normal one, $\rho = s_1$ and the two
statements coincide. For a general matrix only $\rho\leq s_1$ holds, and the bound of
\cref{rnn-bptt:prop:bound} is the one that the argument proves: a matrix whose eigenvalues all
have modulus smaller than $1$ can still multiply a vector by a large factor for many steps
before the product decays (\cref{rnn-bptt:exr:nonnormal}).

\begin{figure}[ht]
\centering
\begin{tikzpicture}
\begin{axis}[width=0.78\textwidth, height=6cm, xlabel={distance $k = T - t$}, ylabel={$\log_{10}\norm{\partial h_T/\partial h_{T-k}}_2$},
  xmin=0, xmax=60, ymin=-12, ymax=3, grid=major, grid style={gray!20},
  legend pos=south west, legend style={font=\scriptsize}, tick label style={font=\scriptsize}, label style={font=\small}]
\addplot[thick, bookpitfall, mark=none] coordinates {(0,0.00) (3,0.59) (6,0.74) (9,0.86) (12,0.94) (15,1.02) (18,1.04) (21,1.12) (24,1.14) (27,1.22) (30,1.27) (33,1.35) (36,1.41) (39,1.44) (42,1.49) (45,1.52) (48,1.58) (51,1.62) (54,1.69) (57,1.77) (60,2.07)};
\addlegendentry{$s_1 = 2$}
\addplot[thick, bookexample, mark=none] coordinates {(0,0.00) (3,-0.05) (6,-0.36) (9,-0.68) (12,-1.07) (15,-1.44) (18,-1.76) (21,-2.09) (24,-2.58) (27,-2.98) (30,-3.41) (33,-3.84) (36,-4.26) (39,-4.67) (42,-5.14) (45,-5.53) (48,-5.92) (51,-6.37) (54,-6.77) (57,-7.18) (60,-7.56)};
\addlegendentry{$s_1 = 1.1$}
\addplot[thick, bookaccent, mark=none] coordinates {(0,0.00) (3,-0.28) (6,-0.77) (9,-1.37) (12,-1.91) (15,-2.47) (18,-3.04) (21,-3.63) (24,-4.28) (27,-4.90) (30,-5.52) (33,-6.14) (36,-6.73) (39,-7.38) (42,-8.02) (45,-8.63) (48,-9.27) (51,-9.89) (54,-10.50) (57,-11.08) (60,-11.64)};
\addlegendentry{$s_1 = 0.9$}
\end{axis}
\end{tikzpicture}
\caption{Norm of the Jacobian product for a $\tanh$ RNN with $m = 50$ units, $W_{hh} = s_1Q$ with $Q$ a random
orthogonal matrix (all singular values equal $s_1$), random scalar inputs, $T = 60$; average of $\log_{10}$ over
$20$ networks. With $s_1 = 0.9$ the decay is guaranteed by \cref{rnn-bptt:prop:bound}. With $s_1 = 1.1$ the
bound allows growth, but saturation ($\phi' < 1$) makes the product vanish anyway. With $s_1 = 2$ it grows.}
\label{rnn-bptt:fig:vanishing}
\end{figure}

\begin{pitfall}[What vanishes]
``The gradient vanishes'' is loose language. The gradient $\nabla_{W_{hh}}L$ is a sum over pairs $(t,s)$
with $t\leq s$ of terms involving $\partial h_s/\partial h_t$. When the Jacobian products decay, the terms with $s - t$
large become negligible, and the sum is dominated by the short-range terms, which can
be perfectly large. The gradient does not vanish; it becomes \emph{blind to long-range
dependencies}. Gradient descent then improves the short-range predictions and has no
signal to improve the long-range ones.
\end{pitfall}

\paragraph{Remedies.} For exploding gradients, \textbf{gradient norm clipping} is simple and
effective (Pascanu et al., 2013): compute the full gradient $g$ (all parameters together) and,
if $\norm g > c$ for a threshold $c$, replace it by $c\,g/\norm g$. The direction is kept and only the step
length is bounded. In PyTorch this is \texttt{torch.nn.utils.clip\_grad\_norm\_(model.parameters(), c)},
called between \texttt{backward()} and \texttt{step()}. Clipping does nothing against vanishing
contributions. For those the architecture must change: the gated networks of
\cref{ch:gated-rnn} add a path through time whose Jacobian is not a product of
$\diag(\phi')\,W_{hh}$.

\section{Truncated BPTT}\label{rnn-bptt:sec:truncated}

For a very long sequence (a whole book, an hour of audio), storing all the states for one
backward pass is impossible, and one update per sequence is too few. \emph{Truncated BPTT}
(Williams and Peng, \emph{An Efficient Gradient-Based Algorithm for On-Line Training
of Recurrent Network Trajectories}, 1990) cuts the sequence into chunks of length $k$ and processes them in order:
\begin{itemize}[nosep]
  \item the forward pass of each chunk starts from the last hidden state of the previous
        chunk, so the state still carries information from the whole past;
  \item the backward pass stops at the beginning of the chunk: the incoming state is
        treated as a constant input;
  \item the parameters are updated after every chunk.
\end{itemize}
The gradient computed for a chunk is therefore the exact gradient of the chunk's losses
with respect to the parameters \emph{used inside the chunk}, with the initial state fixed. It
ignores the effect of the parameters on the initial state, so dependencies longer than $k$
steps receive no gradient. Each chunk costs $O(km)$ memory instead of $O(Tm)$.

\begin{lstlisting}[language=Python]
h = torch.zeros(batch_size, hidden_size)
for chunk_x, chunk_y in chunks:                 # consecutive pieces of the sequence
    h = h.detach()                              # keep the value, cut the graph
    loss = 0.0
    for x_t, y_t in zip(chunk_x, chunk_y):
        h = torch.tanh(x_t @ Wxh.T + h @ Whh.T + bh)
        loss = loss + F.cross_entropy(h @ Who.T + bo, y_t)
    optimizer.zero_grad(); loss.backward()
    torch.nn.utils.clip_grad_norm_(params, max_norm=1.0)
    optimizer.step()
\end{lstlisting}
\texttt{detach()} returns a tensor with the same value and no history: the forward
computation continues, but the backward pass of the next chunk stops there. Without it, the
second \texttt{backward()} would try to go back through the graph of the first chunk, which
has already been freed. (The rows of a mini-batch are vectors, hence $x_tW_{xh}^\top$ instead of
$W_{xh}x_t$, as in \cref{backpropagation:sec:batch}.)

\section{Appendix: a proof of the chain rule}\label{rnn-bptt:sec:chain-proof}

Everything in this chapter rests on the chain rule (\cref{backpropagation:thm:chain}). We
prove it, because the most common proof contains a gap.

\begin{theorem}\label{rnn-bptt:thm:chain}
Let $f:\R^n\to\R^m$ be differentiable at $x$ and $g:\R^m\to\R^p$ differentiable at $y = f(x)$, with Jacobians $A = J_f(x)$
and $B = J_g(y)$. Then $g\circ f$ is differentiable at $x$ and $J_{g\circ f}(x) = BA$.
\end{theorem}

\begin{proof}
Differentiability of $f$ at $x$ means $f(x + v) = f(x) + Av + r(v)$ with $\norm{r(v)}/\norm v\to0$ as $v\to0$. For $g$, define
\[
  \varepsilon(k) = \frac{g(y + k) - g(y) - Bk}{\norm k}\quad(k\neq0),\qquad\varepsilon(0) = 0,
\]
so that $g(y + k) = g(y) + Bk + \norm k\,\varepsilon(k)$ holds for \emph{every} $k$, including $k = 0$, and $\varepsilon(k)\to0 = \varepsilon(0)$
as $k\to0$: $\varepsilon$ is continuous at $0$. Put $k(v) = f(x + v) - f(x) = Av + r(v)$. Then
\[
  g(f(x + v)) = g(y + k(v)) = g(y) + BAv + \underbrace{Br(v) + \norm{k(v)}\,\varepsilon(k(v))}_{R(v)} ,
\]
and it remains to show $\norm{R(v)}/\norm v\to0$. First, $\norm{Br(v)}/\norm v\leq\norm B\,\norm{r(v)}/\norm v\to0$. Second,
\[
  \frac{\norm{k(v)}}{\norm v}\leq\norm A + \frac{\norm{r(v)}}{\norm v},
\]
which is \emph{bounded} for small $v$ (it need not converge), and $k(v)\to0$ as $v\to0$, so $\varepsilon(k(v))\to0$ by
continuity of $\varepsilon$ at $0$. A bounded quantity times one that tends to $0$ tends to $0$.
\end{proof}

The gap of the usual proof is the division by $\norm{k(v)}$: one writes the remainder of $g$ as
$\norm k\,\varepsilon(k)$ only for $k\neq0$, and $k(v) = f(x + v) - f(x)$ can be $0$ for $v\neq0$ arbitrarily close to $0$ (take $f$
constant). Defining $\varepsilon(0) = 0$ removes the problem. A second, related mistake is to claim
that $\norm{k(v)}/\norm v$ tends to $\norm{Av}/\norm v$ or to $\norm A$: it does not converge in general (it depends on the
direction of $v$), and only its boundedness is needed.

\begin{chaptersummary}
\begin{itemize}
  \item An RNN applies the same affine map and activation at every step,
        $h_t = \phi(W_{hh}h_{t-1} + W_{xh}x_t + b_h)$; $h_t$ is a fixed-size summary of $x_{1:t}$ and the number of parameters
        does not depend on the length.
  \item Unrolled, the RNN is a computational graph with shared weights. Backpropagation on it
        (BPTT) gives $\bar h_t = W_{ho}^\top\bar o_t + W_{hh}^\top\delta_{t+1}$, $\delta_t = \bar h_t\odot\phi'(a_t)$, and parameter gradients that are
        sums over time, such as $\nabla_{W_{hh}}L = \sum_t\delta_th_{t-1}^\top$.
  \item $\bar h_t$ is a \emph{total} derivative; the loss at time $s$ reaches $h_t$ through the product of
        $s - t$ Jacobians $\diag(\phi'(a_r))W_{hh}$.
  \item $\norm{\partial h_s/\partial h_t}_2\leq(\gamma s_1)^{s-t}$: if $\gamma s_1 < 1$ the long-range contributions vanish
        geometrically; $\gamma s_1 > 1$ is necessary, not sufficient, for explosion. Singular values, not
        eigenvalues, control the bound.
  \item Gradient norm clipping bounds exploding steps; truncated BPTT bounds memory and
        ignores dependencies longer than the chunk; vanishing contributions call for gated
        architectures.
\end{itemize}
\end{chaptersummary}

\section*{Exercises}
\addcontentsline{toc}{section}{Exercises}

\begin{exercise}\label{rnn-bptt:exr:1}
An RNN reads sequences of $100$-dimensional word vectors with $m = 128$ hidden units and
predicts one of $K = 10$ tags at every step. Count its parameters. How many would an MLP with
one hidden layer of $128$ units need to read sentences of $50$ words at once and output
$50$ tags?
\solutionlink{rnn-bptt:exr:1}
\end{exercise}

\begin{exercise}\label{rnn-bptt:exr:2}
In a many-to-one RNN, the loss is $L = \ell_T$ only. Write the BPTT recursion and the
gradient of $W_{ho}$. Express $\bar h_t$ with the Jacobian products.
\solutionlink{rnn-bptt:exr:2}
\end{exercise}

\begin{exercise}\label{rnn-bptt:exr:3}
In \cref{rnn-bptt:ex:numbers}, recompute $\bar h_3$ and $\delta_3$ from $\hat y_3$ and $h_3$, and check the
contribution $\delta_3h_2^\top$.
\solutionlink{rnn-bptt:exr:3}
\end{exercise}

\begin{exercise}\label{rnn-bptt:exr:nonnormal}
Let $W = \begin{bmatrix}0.5 & 3\\ 0 & 0.5\end{bmatrix}$ and consider the linear RNN $h_t = Wh_{t-1}$ (so $\partial h_s/\partial h_t = W^{s-t}$). Give the
eigenvalues of $W$, prove that $W^k = \begin{bmatrix}0.5^k & 6k\,0.5^k\\ 0 & 0.5^k\end{bmatrix}$, and compute the largest entry of $W^k$ for
$k = 1,\dots,5$. What does this say about the eigenvalue condition for vanishing gradients?
\hint{Write $W = 0.5(I + N)$ with $N^2 = 0$.}
\solutionlink{rnn-bptt:exr:nonnormal}
\end{exercise}

\begin{exercise}\label{rnn-bptt:exr:5}
Truncated BPTT with chunks of length $k = 1$. Show that the computed gradient of $W_{hh}$ is
$\sum_t\bigl((W_{ho}^\top\bar o_t)\odot\phi'(a_t)\bigr)h_{t-1}^\top$. Which term of \cref{rnn-bptt:prop:unrolled} is kept?
\solutionlink{rnn-bptt:exr:5}
\end{exercise}

\begin{exercise}\label{rnn-bptt:exr:6}
Let $f:\R\to\R^2$, $f(v) = (v^2\sin(1/v),\,0)$ for $v\neq0$, $f(0) = 0$. Show that $f$ is differentiable at $0$ with $J_f(0) = 0$, that $k(v) = f(v) - f(0)$
vanishes for $v = 1/(j\pi)$, $j\in\N^*$, and explain why this example defeats a proof of the chain rule that
divides by $\norm{k(v)}$.
\solutionlink{rnn-bptt:exr:6}
\end{exercise}
\end{notechapter}
